\section{Recomendaciones}

A partir de los resultados obtenidos y de las limitaciones identificadas,
se plantean las siguientes recomendaciones:

\begin{enumerate}
    \item Completar la verificación del flujo radiográfico en producción
    después de publicar la revisión con las últimas correcciones.
    Esta comprobación debe incluir la creación de la solicitud,
    la generación del enlace o código QR, la entrega directa a Supabase
    Storage y la consulta del archivo por un odontólogo autorizado.
    Se recomienda conservar evidencias de los recorridos correctos
    y de los rechazos por vencimiento, reutilización y falta de permisos.

    \item Consolidar las funcionalidades RF-07 y RF-08, clasificadas
    como \textit{Could have}, correspondientes a la trazabilidad
    del proceso radiográfico y la consulta del estado de solicitudes.
    Puesto que estas funciones ya están presentes en la versión web,
    las mejoras posteriores pueden orientarse a reportes exportables
    e indicadores de solicitudes pendientes y completadas, previa
    validación de su utilidad con el personal de la clínica.

    \item Delimitar el componente móvil a las necesidades de consulta
    del odontólogo, especialmente los pacientes asignados y sus
    radiografías. Su evaluación debe comprobar el uso de la misma API,
    la autorización aplicada en el servidor y el cierre de sesión,
    además de documentar la distribución y actualización del instalable.

    \item Evaluar la escalabilidad mediante pruebas de carga con
    usuarios concurrentes y un volumen representativo de registros.
    A partir de sus resultados, revisar las consultas, los índices,
    la paginación y el uso de conexiones a PostgreSQL. También se
    recomienda supervisar los errores, los tiempos de respuesta
    y el crecimiento del almacenamiento, porque las diez mediciones
    realizadas no permiten estimar el comportamiento bajo alta demanda.

    \item Realizar una adopción gradual en la Clínica Odontológica
    San Martín, acompañada de capacitación para administradores,
    odontólogos y recepción, así como de instrucciones de entrega
    para CERPAX. Se recomienda definir responsables de soporte,
    revisar periódicamente los permisos y establecer procedimientos
    de respaldo y restauración de los registros y archivos.
    Durante el uso piloto deben recogerse observaciones del personal
    y métricas de operación para orientar las mejoras posteriores.
\end{enumerate}
